\section{Análisis de Requerimientos}

En esta sección se describen los requerimientos funcionales y no funcionales del sistema propuesto, los cuales establecen las capacidades, restricciones y criterios de desempeño que debe cumplir la solución desarrollada. Estos requerimientos sirven como base para el diseño, implementación y evaluación del sistema, asegurando la trazabilidad entre las necesidades del evaluador y la arquitectura final.

% =========================================================
\subsection{Requerimientos Funcionales}
% =========================================================

Los requerimientos funcionales definen las capacidades y procesos que el sistema debe implementar para cumplir con las funciones de análisis, gestión de información y evaluación del código en el contexto académico.

\begin{itemize}
    \item \textbf{RF-01:} El sistema debe permitir el registro e inicio de sesión de docentes mediante credenciales cifradas para garantizar un espacio de trabajo privado.
    
    \item \textbf{RF-02:} El sistema debe permitir al docente la creación y gestión de un listado de ejercicios o problemas de programación (Repositorio Personal).
    
    \item \textbf{RF-03:} El sistema debe permitir la carga de código fuente en formato .c, .cpp y .txt asociados a un ejercicio específico.
    
    \item \textbf{RF-04:} El sistema debe realizar la normalización adecuada del código, eliminando comentarios para la extracción de grafos y preservando las estructuras de estilo para el análisis de autoría.
    
    \item \textbf{RF-05:} El sistema debe generar representaciones basadas en Grafos de Flujo de Datos (DFG) y n-gramas de caracteres.
    
    \item \textbf{RF-06:} El sistema debe realizar el análisis de similitud semántica mediante el modelo GraphCodeBERT para la generación de vectores de características (embeddings), cuantificando la proximidad entre ellos utilizando la métrica de similitud del coseno.
    
    \item \textbf{RF-07:} El sistema debe realizar el análisis estilométrico mediante el modelo CharCNN para calcular la probabilidad de autoría.
    
    \item \textbf{RF-08:} El sistema debe solicitar y recibir múltiples versiones funcionalmente equivalentes del problema proporcionado por el evaluador a través de un modelo LLM externo (Trayectoria A1).
    
    \item \textbf{RF-09:} El sistema debe permitir el almacenamiento persistente de los embeddings y resultados en una base de datos vectorial junto con sus metadatos lógicos.
    
    \item \textbf{RF-10:} El sistema debe generar un reporte bimodal que integre el índice de similitud semántica y la probabilidad de autoría de forma visual.
    
    \item \textbf{RF-11:} El sistema debe emitir alertas de integridad de forma automática cuando se detecten discrepancias críticas entre la lógica y el estilo de la entrega (Trayectorias A4 y A5).
    
    \item \textbf{RF-12:} El sistema debe permitir la consulta histórica de reportes de análisis realizados previamente sin necesidad de re-procesar los archivos fuente.
\end{itemize}

% =========================================================
\subsection{Requerimientos No Funcionales}
% =========================================================

Los requerimientos no funcionales establecen las restricciones, atributos de calidad y condiciones de desempeño que debe cumplir el sistema propuesto durante su operación.

\begin{itemize}
    \item \textbf{RNF-01:} Seguridad de datos: El sistema debe cifrar las contraseñas de los usuarios y utilizar tokens de sesión (JWT) para proteger la comunicación entre el cliente y el servidor.
    
    \item \textbf{RNF-02:} Persistencia híbrida: El sistema debe contar con un esquema de base de datos dual (SQL para metadatos de usuario y Vectorial para embeddings de código).
    
    \item \textbf{RNF-03:} El modelo GraphCodeBERT debe alcanzar un F1-Score mayor a 0.85 en tareas de detección de clones.
    
    \item \textbf{RNF-04:} El modelo CharCNN debe alcanzar una precisión de autoría mayor al 0.8.
    
    \item \textbf{RNF-05:} El sistema debe operar bajo los estándares C11 y C++17 para garantizar compatibilidad con el currículo de ESCOM.
    
    \item \textbf{RNF-06:} El sistema debe mantener una tasa de falsos positivos menor al 10\% mediante el orquestador de métricas.
    
    \item \textbf{RNF-07:} El tiempo de procesamiento total por entrega no debe exceder los 20 segundos.
    
    \item \textbf{RNF-08:} El sistema debe ser compatible con navegadores basados en Chromium versión 110 o superior.
    
    \item \textbf{RNF-09:} La interfaz debe ser responsiva y mostrar indicadores de progreso en tiempo real durante la inferencia de los modelos.
    
    \item \textbf{RNF-10:} El sistema debe garantizar una disponibilidad del 99\% para su uso en periodos de evaluación.
    
    \item \textbf{RNF-11:} El sistema debe validar la integridad y tipo de los archivos en el cliente antes de ser enviados al backend.
\end{itemize}

% =========================================================
\subsection{Reglas de Negocio}
% =========================================================

Las reglas de negocio definen los criterios que guían el funcionamiento y uso ético del sistema dentro del entorno académico.

\begin{itemize}
    \item \textbf{RN-01:} Confidencialidad de los datos: Un docente solo podrá visualizar, editar y consultar los ejercicios y reportes generados bajo su propia cuenta de usuario.
    
    \item \textbf{RN-02:} El sistema se limita a programas de cursos de programación de ESCOM-IPN que sigan los estándares C11 y C++17.
    
    \item \textbf{RN-03:} Integridad de la referencia: Toda entrega debe ser comparada contra al menos una solución de referencia proporcionada por el docente para ser considerada válida.
    
    \item \textbf{RN-04:} La decisión final de evaluación es responsabilidad exclusiva del docente; el sistema actúa únicamente como una herramienta de apoyo a la auditoría académica.
    
    \item \textbf{RN-05:} Se permite el uso de múltiples soluciones generadas por LLM para ampliar el espectro de comparación lógica, siempre que el docente las valide inicialmente.
    
    \item \textbf{RN-06:} La autoría del código se infiere mediante la desviación del vector estilométrico respecto al estándar humano o sintético identificado por el canal B.
\end{itemize}